\section{模型的评价}

\subsection{模型的优点}

\begin{itemize}
	\item \textbf{约束刻画完备}：将土地容量、季节适应性、水浇地二季复种、重茬回避、豆类三年轮作、销售上限等农业制度约束全部显式纳入0-1整数规划，方案可直接落地执行。
	\item \textbf{统一建模框架}：问题1、2、3共用同一套种植决策变量与核心约束，仅改变目标函数与参数口径（确定性$\to$随机$\to$相关随机），便于横向对比与结果解释。
	\item \textbf{风险量化科学}：采用CVaR（条件风险价值）度量尾部损失，并给出其线性化表达式\cite{cvar2000,cvar2002}，可在大型整数规划中直接求解，兼顾期望收益与最坏情景，避免了单纯方差度量忽视尾部风险的缺陷。
	\item \textbf{相关结构可解释}：问题3从替代、互补、产销联动等经济学机制构造相关矩阵，并用Cholesky分解生成相关情景，机制明确、可解释性强。
	\item \textbf{求解质量有保证}：采用商用级开源求解器 HiGHS 分支定界，问题1收敛到$1\%$相对间隙，问题2、3在$4\%$间隙内求解，配合约束程序化校验，结果可信。
\end{itemize}

\subsection{模型的缺点}

\begin{itemize}
	\item \textbf{分布假设简化}：题目仅给波动区间，本文统一采用截断标准正态抽样（为保证问题2、3边际一致），真实分布可能偏离该假设，导致风险估计存在偏差。
	\item \textbf{情景规模受限}：受求解规模限制，问题2、3取$N=30$个情景，对联合分布尾部的刻画有限，风险统计量存在约$\pm15\%$量级的抽样误差；情景数增大时求解时间快速上升。
	\item \textbf{相关结构静态}：问题3的相关矩阵为人工设定的静态参数，未利用历史产量-销量数据估计，相关强度的实证依据不足。
	\item \textbf{单作物满块假设}：实际可部分种植、间作套种，本文为控制规模做了简化。
\end{itemize}

\section{模型的改进与推广}

\subsection{模型的改进}

未来可：(i) 用历史多年（2020--2023）产量与销售数据估计相关系数与分布参数，替代人工设定；(ii) 采用Benders分解、列生成等算法突破情景规模限制，将$N$提升至数百个情景以获得更稳的尾部估计；(iii) 引入多阶段（逐年滚动）随机规划，允许方案随年度信息更新而调整；(iv) 将单作物满块约束松弛为部分种植，扩展为连续面积变量与整数变量的混合模型。

\subsection{模型的推广}

本文框架具有较强通用性：更换作物与地块数据即可用于其他地区的种植结构优化；增加劳动力、水资源、生态补偿等维度即可适应更复杂的农业生产决策；CVaR+随机规划的框架还可推广到供应链库存、能源调度、金融投资组合等一切"现在决策、事后实现"的资源配置问题，具有广阔的推广前景。
